\documentclass[11pt,a4paper]{article}

% --- ENCODING & LANG ---
\usepackage[utf8]{inputenc}
\usepackage[T1]{fontenc}
\usepackage[ngerman,main=english]{babel}
\usepackage{lmodern}
\usepackage{microtype}

% --- MATH & CORE UTILITIES ---
\usepackage{amsmath,amssymb,amsfonts,amsthm}
\usepackage{cite}
\usepackage{hyperref}
\usepackage{enumitem}

% --- THEOREM ENVIRONMENTS ---
\newtheorem{theorem}{Theorem}[section]
\newtheorem{definition}[theorem]{Definition}
\newtheorem{lemma}[theorem]{Lemma}
\newtheorem{remark}[theorem]{Remark}
\newtheorem{example}[theorem]{Example}

\title{%
  \textbf{Nicht-Metaphysics: Paper 0 -- Two Modes of Consistency}\\
  \large A Formal Framework for Frame-Consistency and Inversion Stability%
}

\author{%
  \textbf{xerx593} \quad \& \quad \textbf{Non-Human Interlocutors}%
}
\IfFileExists{release_dates.def}{% ======================================================================
% Central Production Release Date Registry for nicht-physics
% Maps root filename (\jobname) to first_release_date & last_release_date
% ======================================================================

% Helper macro to set dates per file stem
\newcommand{\registerarticledates}[3]{%
  \expandafter\def\csname firstdate@#1\endcsname{#2}%
  \expandafter\def\csname lastdate@#1\endcsname{#3}%
}

% --- Production Metadata Registry ---
\registerarticledates{article_forces_and_constants}{August 24, 2026}{August 27, 2026}
\registerarticledates{article_yang_mills_mass_gap}{August 24, 2026}{August 24, 2026}
\registerarticledates{article_apophatic_optimization}{August 25, 2026}{August 25, 2026}
\registerarticledates{article_apophatic_dynamics}{August 27, 2026}{August 27, 2026}
\registerarticledates{article_triade_intractabilities}{August 27, 2026}{August 27, 2026}
\registerarticledates{article_meth_space}{August 29, 2026}{August 29, 2026}
\registerarticledates{article_meth_riemann}{August 29, 2026}{August 29, 2026}
\registerarticledates{article_meth_bsd}{August 29, 2026}{August 29, 2026}
\registerarticledates{article_meth_hodgepoincare}{August 29, 2026}{August 29, 2026}

% --- Automatic Resolution via \jobname (with Apophatic Fallback) ---
\expandafter\ifx\csname firstdate@\jobname\endcsname\relax
  % Fallback for unregistered files or draft builds
  \newcommand{\firstpublishdate}{Some day A.D.}
  \newcommand{\apophaticdate}{\firstpublishdate}
\else
  % Load registered production dates
  \edef\firstpublishdate{\csname firstdate@\jobname\endcsname}
  \edef\apophaticdate{\csname lastdate@\jobname\endcsname}
\fi}{}

\date{\apophaticdate}

\begin{document}

\maketitle

\begin{abstract}
This paper introduces a formal distinction between two modes of structural 
evaluation in formal systems. \emph{Frame-Consistency} ($\mathrm{C}_{\mathrm{frame}}$) 
denotes standard internal consistency ($S \nvdash \bot$) within a fixed 
axiomatic frame $S$. \emph{Inversion Stability} ($\mathrm{C}_{\mathrm{inv}}$) evaluates 
whether a system remains closed and non-contradictory under a specified 
domain-specific counter-operation or inversion $I$. We define the 
\emph{Invariant Baseline} $B_0$ as a sub-structure preserved as a set 
under $I$ ($I(B_0) = B_0$). We show that non-closure under an inversion 
operator is distinct from internal syntactic inconsistency. This distinction 
provides a domain-agnostic foundation for domain-specific investigations 
in number theory, algebraic geometry, and physics.
\end{abstract}

\vspace{0.3cm}

\begin{otherlanguage}{ngerman}
\begin{abstract}
Dieses Paper führt eine formale Unterscheidung zwischen zwei Moden der 
strukturellen Evaluierung in formalen Systemen ein. \emph{Rahmen-Konsistenz} 
($\mathrm{C}_{\mathrm{frame}}$) bezeichnet die gewöhnliche innere Konsistenz ($S \nvdash \bot$) 
innerhalb eines festgelegten axiomatischen Rahmens $S$. 
\emph{Inversions-Stabilität} ($\mathrm{C}_{\mathrm{inv}}$) prüft, ob ein System unter einer 
spezifizierten bereichsspezifischen Gegenoperation oder Inversion $I$ 
abgeschlossen und widerspruchsfrei bleibt. Wir definieren die 
\emph{Invariante Basis} $B_0$ als eine Unterstruktur, die unter $I$ als 
Menge invariant bleibt ($I(B_0) = B_0$). Wir zeigen, dass 
Nicht-Abgeschlossenheit unter einem Inversionsoperator nicht mit innerer 
syntaktischer Inkonsistenz identisch ist. Diese Unterscheidung
bietet eine bereichsunabhängige Grundlage für bereichsspezifische Untersuchungen
in der Zahlentheorie, der algebraischen Geometrie und Physik.
\end{abstract}
\end{otherlanguage}

\clearpage

% ============================================================
\section{Introduction}
% ============================================================

Standard mathematical reasoning begins by fixing a formal system 
$S = (\mathcal{L}, \mathcal{A}, \mathcal{R})$ with language $\mathcal{L}$, 
axioms $\mathcal{A}$, and derivation rules $\mathcal{R}$. Consistency is 
evaluated internally within the theoretical deductive closure $\text{Th}(S)$.

However, formal frameworks frequently generate intrinsic operations---such 
as additive inversion, multiplicative duality, or functional 
reflections---whose outputs may extend beyond the domain of $S$. Internal 
consistency within $S$ does not guarantee structural closure under such 
counter-operations.

We formalize this distinction via two operational perspectives:
\[
\boxed{
\text{Frame-Consistency } (\mathrm{C}_{\mathrm{frame}})
\quad\neq\quad
\text{Inversion Stability / Closure } (\mathrm{C}_{\mathrm{inv}})
}
\]

% ============================================================
\section{Two Modes of Consistency}
% ============================================================

Let $S = (\mathcal{L}, \mathcal{A}, \mathcal{R})$ be a formal system, 
and let $\vdash_S$ denote syntactic provability. Let $\bot$ denote logical 
absurdity.

\subsection{Frame-Consistency ($\mathrm{C}_{\mathrm{frame}}$)}

\begin{definition}[Frame-Consistency]
A formal system $S$ is \emph{Frame-Consistent} (denoted $\mathrm{C}_{\mathrm{frame}}$) if 
no contradiction is derivable within its operational rules:
\[
S \nvdash_S \bot.
\]
\end{definition}

\subsection{Inversion Stability ($\mathrm{C}_{\mathrm{inv}}$)}

Let $I: \mathcal{D} \to \mathcal{D}'$ be a domain-specific inversion 
operator acting on a domain $\mathcal{D} \subseteq \mathcal{L}$.

\begin{definition}[Invariant Baseline $B_0$]
An \emph{Invariant Baseline} $B_0 \subseteq \mathcal{D}$ for an inversion 
operator $I$ is a sub-structure or domain that is invariant as a set 
under $I$:
\[
I(B_0) = B_0.
\]
\end{definition}

\begin{remark}
An invariant baseline $B_0$ does not require pointwise fixed points 
($\forall x \in B_0, I(x) = x$). It requires set-wise invariance under 
the action of $I$.
\end{remark}

\begin{definition}[Inversion Extension and $\mathrm{C}_{\mathrm{inv}}$]
Let $S_I$ denote the minimal inductive extension of $S$ closed under the 
action of $I$. System $S$ is \emph{Inversion-Stable} with respect to $I$ 
(denoted $\mathrm{C}_{\mathrm{inv}}$) if:
\begin{enumerate}[label=(\roman*)]
    \item $S$ is closed under $I$ (i.e., $I(\mathcal{D}) \subseteq \mathcal{D}$), and
    \item The extended theory $S_I$ is Frame-Consistent ($S_I \nvdash \bot$).
\end{enumerate}
\end{definition}

% ============================================================
\section{Non-Closure vs. Inconsistency}
% ============================================================

A primary source of conceptual confusion in structural logic is the 
conflation of domain non-closure with internal logical inconsistency.

\begin{example}[Arithmetic Non-Closure]
Consider the system $S = (\mathbb{N}, +)$. $S$ is $\mathrm{C}_{\mathrm{frame}}$ (internally 
consistent). Define the additive inversion operator $I(a, b) = a - b$.
For $a < b$, $I(a,b) \notin \mathbb{N}$. The system $S$ is not closed 
under $I$, making it unclosed under inversion. However, this non-closure 
does not generate logical absurdity ($\bot$) within $S$; it merely 
defines a boundary prompting extension to $\mathbb{Z}$.
\end{example}

\begin{theorem}[Separation of $\mathrm{C}_{\mathrm{frame}}$ and $\mathrm{C}_{\mathrm{inv}}$]
Frame-consistency ($\mathrm{C}_{\mathrm{frame}}$) does not imply Inversion Stability ($\mathrm{C}_{\mathrm{inv}}$).
\end{theorem}

\begin{proof}
Let $S = (\mathbb{N}, +)$ and $I(a, b) = a - b$. $S \nvdash \bot$, so $S$ is 
$\mathrm{C}_{\mathrm{frame}}$. However, $I(2, 3) = -1 \notin \mathbb{N}$, so $S$ is not closed 
under $I$. Thus, $S$ is $\mathrm{C}_{\mathrm{frame}}$ but not $\mathrm{C}_{\mathrm{inv}}$.
\end{proof}

% ============================================================
\section{Core Reference Mapping}
% ============================================================

This framework allows downstream domain-specific papers to evaluate stability 
under domain-appropriate inversion operators $I$ and invariant baselines $B_0$:

\begin{table}[h!]
\centering
\small
\begin{tabular}{lll}
\hline
\textbf{Domain} & \textbf{Inversion Operator $I$} & \textbf{Invariant Baseline $B_0$} \\ \hline
Analytic Number Theory & Reflection $s \mapsto 1-s$ & Critical Line $\operatorname{Re}(s) = 1/2$ \\
Elliptic Curves & Analytic Shift $s \mapsto 2-s$ & Central Point $s = 1$ \\
Vector Spaces & Dual Map $V \mapsto V^*$ & Self-Dual Subspaces $V \cong V^*$ \\ \hline
\end{tabular}
\end{table}

% ============================================================
\section{Conclusion}
% ============================================================

Paper 0 establishes a precise distinction between internal syntactic rule 
compliance ($\mathrm{C}_{\mathrm{frame}}$) and structural closure under counter-operations 
($\mathrm{C}_{\mathrm{inv}}$). By defining the invariant baseline as an invariant set 
$I(B_0) = B_0$, this framework provides a solid reference for domain-specific 
applications without requiring metaphysical postulates or categorical 
over-extensions.

\clearpage
\bibliographystyle{unsrt}
\bibliography{references}

\end{document}
